\section{Boussinot's Reactive Control Model}
\label{sec:boussinot-control}
Boussinot's reactive control model organizes concurrency around logical
instants. Execution proceeds as a sequence of reaction cycles: within an
instant, all runnable activities execute until they suspend or the system
quiesces, and only after that quiescence can absence be observed. This notion
of instant provides a global time base that makes reactions comparable and
deterministic at instant boundaries \cite{fairthreads2003}.

\paragraph{Signals (unvalued)}
To simplify the presentation, we consider unvalued signals. A signal is either
present or absent for the entire duration of an instant. When a signal is
emitted, it becomes present for all components; when it is not emitted, it
remains absent for all components. The effective order of emissions within the
instant must not affect the observable behavior: presence is a global property
for the instant, not a per-thread event. This makes synchronization explicit
and avoids order-sensitive communication.

\paragraph{Suspension}
Threads are cooperative: they execute until they reach a suspension point and
then yield control. Suspension can be explicit (e.g., waiting for a signal or
for the next instant). The key property is that suspension preserves the
instant structure: a suspended thread resumes only at a well-defined instant
boundary or when a relevant signal becomes present.

\paragraph{Algorithmic view (abstract)}
We use a small algorithmic notation that mirrors Boussinot's control vocabulary.
Behaviors (threads) are composed from a core set of control operators. The
notation hides values when they are not essential to control and uses abstract
behaviors inside control structures:
\begin{lstlisting}[style=tempo]
(* Abstract control language; B is an abstract behavior. *)
B ::= skip
    | emit s          (* mark signal s present in current instant *)
    | await s         (* resume at next instant where s is present *)
    | pause           (* resume at the next instant *)
    | B1; B2          (* sequence *)
    | parallel [B1; ...; Bn]
    | when s do B     (* guard: block until s is present *)
    | watch s do B    (* weak preemption: stop B after this instant if s emits *)
\end{lstlisting}

The constructs are interpreted at instant boundaries. \texttt{emit s} makes
\texttt{s} present for the remainder of the instant. \texttt{await s} suspends
until the next instant where \texttt{s} is present, while \texttt{pause} always
advances to the next instant. In \texttt{when s do B}, the guard is a blocking
condition: if \texttt{s} is absent, the entire behavior suspends and can only
resume at a future instant where \texttt{s} is present. In \texttt{watch s do B},
the body \texttt{B} runs normally in the current instant, but if \texttt{s} is
emitted during that instant, \texttt{B} is not allowed to continue into the next
instant (weak preemption).

This abstract view lets us express control patterns without committing to a
concrete evaluation order. For example, \texttt{when s do \{B1; B2\}} denotes an
entire guarded behavior, and \texttt{watch s do \{parallel [B1; B2]\}} denotes a
preemptible concurrent block.

\paragraph{Micro-examples (abstract)}
\begin{lstlisting}[style=tempo]
(* Suspension across instants *)
emit s; pause; await s

(* Guarded behavior: runs only if s is present *)
when s do (emit out)

(* Weak preemption: body stops after instant if k is emitted *)
watch k do (emit tick; pause; emit tick)

(* Disjunctive enabling via parallel guards *)
parallel [
  (when g1 do B);
  (when g2 do B)
]
\end{lstlisting}

\paragraph{Preemption and Esterel}
Control languages such as Esterel rely on preemption to enforce control
structure. In Boussinot's setting, preemption is \emph{weak}: if a preempting
signal is emitted in instant $t$, a preempted computation may complete its
current instant but must not continue into $t{+}1$. The alternative would be
\emph{strong} preemption, where a computation is stopped immediately within
the same instant. Weak preemption is the appropriate choice here because the
model treats the instant as an atomic reaction: all threads share the same
present/absent view for the whole instant, and mid-instant interruption would
reintroduce order dependence and break determinism \cite{esterel1992,fairthreads2003}.

This model yields a control-centric view of concurrency: threads cooperate,
rendezvous on signal presence, and advance in lockstep from one instant to the
next. \tempo{} adopts this discipline as the semantic baseline for its runtime.

\paragraph{From FairThreads to ReactiveML (abstract view)}
ReactiveML lifts this control model into a higher-order ML setting while
preserving the delayed-absence discipline. It exposes a small set of reactive
operators---\texttt{signal}, \texttt{emit}, \texttt{await}, \texttt{pause},
\texttt{when}, and \texttt{watch}---that directly encode the instant-based
control structure. These constructs compile to a runtime that enforces the
instant discipline and yields deterministic behavior at instant boundaries,
while allowing higher-order composition over reactive processes
\cite{reactiveml2005}.

At the same level of abstraction, this places ReactiveML in the control-flow
lineage (Esterel/FairThreads), in contrast to dataflow approaches such as FRP
that emphasize time-varying values and stream composition \cite{frp1997}.
Recent FRP work has refined operational guarantees such as causality and space
usage (e.g., \cite{perez2016frp,bahr2019ratt,bahr2022modalfrp,bahr2023async}),
but these advances are orthogonal to the control-centric lineage considered
here.

ReactiveML also benefits from language-level integration: its compiler can
enforce discipline constraints and optimize the scheduling structure
statically. \tempo{} chooses a different point in the design space by reusing
OCaml 5.3 algebraic effects to recover similar control constructs as a
library. This enables straightforward embedding in existing OCaml codebases
while leaving stronger static guarantees to future work.
